\begin{claim}{D.2.47}{Proved}{The integral half-class has a nonzero mixed differential}
Put $B=K(\FF_2,2)$, $Y=\Sigma^2B$, $Z=Y\langle5\rangle$,
and $L=Y\langle7\rangle$. Write $U_4=\Sigma^2u_2\in H^4(Y;\FF_2)$.
The actual six-type and its loop space are
\[
\begin{aligned}
 P_6Y&\simeq\fib\bigl(\operatorname{Sq}^3:
       K(\FF_2,4)\longrightarrow K(\FF_2,7)\bigr),
 \\
 H=\Omega P_6Y&\simeq\fib\bigl(v_3^2:
       K(\FF_2,3)\longrightarrow K(\FF_2,6)\bigr).
\end{aligned}
 \tag{D.2.47a}
\]
The fundamental class $a_7\in H^7(L;\ZZ/4)$ of D.2.45 is the
restriction of the actual secondary Pontryagin class on $Z$
obtained from a nullhomotopy of $\mathfrak P_2(U_4)$.
This restriction is independent of that nullhomotopy, up to the
chosen sign of the fundamental generator.

Let $b_8=\beta_4a_7\in H^8(L;\ZZ)$, where $\beta_4$ is the
integral coefficient connecting map, and put
$V_4=\beta_{\ZZ}v_3\in H^4(K(\FF_2,3);\ZZ)$.
For the actual fibration $H\longrightarrow L\longrightarrow Y$,
there is a generator $z_8$ such that
\[
 H^8(H;\ZZ)=\ZZ/4\{z_8\},\qquad
 b_8|_H=V_4^2=2z_8\ne0,\qquad
 (b_8^2)|_H=0.
 \tag{D.2.47b}
\]
Its integral Serre spectral sequence has the nonzero mixed
differential
\[
 d_5(z_8)=\operatorname{Sq}^1U_4\otimes V_4
       \in E_5^{5,4}\cong\FF_2.
 \tag{D.2.47c}
\]
Here the base factor has mod-two coefficients because the
fiber coefficient group $H^4(H;\ZZ)$ is $\ZZ/2$.
Thus the half $z_8$ does not extend to $L$, although its
double does. The value of $b_8^2$ on $L$, and the stable-spherical
pairing in D.2.45, remain undecided by this calculation.
\end{claim}
\begin{proof}
\textbf{The six-type and the actual secondary class.}
The already established low groups are
$\pi_4Y=\FF_2$, $\pi_5Y=0$, and $\pi_6Y=\FF_2$.
The possible invariant lies in
\[
 H^7(K(\FF_2,4);\FF_2)
 =\FF_2\{\operatorname{Sq}^3u_4,
                    \operatorname{Sq}^2\operatorname{Sq}^1u_4\}.
 \tag{D.2.47d}
\]
It is nonzero: otherwise the six-type would be a product,
with two independent degree-six cohomology classes, whereas
$H^6Y$ has dimension one. The map $Y\to P_6Y$ is a homology
isomorphism through degree six. Moreover
\[
 \operatorname{Sq}^3U_4=0,
 \qquad
 \operatorname{Sq}^2\operatorname{Sq}^1U_4
      =\Sigma^2(\operatorname{Sq}^2\operatorname{Sq}^1u_2)\ne0.
\]
This identifies the invariant and proves the first equivalence
in \textup{(D.2.47a)}. Looping gives the second, since
$\operatorname{Sq}^3v_3=v_3^2$.
Equivalently, $P_6Y$ is the pullback of
$K(\ZZ/4,6)\to K(\FF_2,6)$ along $\operatorname{Sq}^2u_4$;
the coefficient connecting map is $\operatorname{Sq}^1$.

Thomas's suspension formula and its second-suspension vanishing give
\[
 S(\mathfrak P_2)=p,\qquad S(p)=0,
 \qquad\mathfrak P_2(U_4)=0
 \tag{D.2.47e}
\]
\cite[Theorem~I, pp.~897--898, and Proposition~6.9,
p.~908]{Thomas1959Pontrjagin}.
Choose an actual nullhomotopy, equivalently a lift of $Y$ to
the fiber of $\mathfrak P_2:K(\FF_2,4)\to K(\ZZ/4,8)$.
Passing to fibers over $K(\FF_2,4)$ gives a map
$a_Z:Z\to K(\ZZ/4,7)$.

To check its coefficient, use the generator
$\alpha=u_4\nu_4:S^7\to Y$, which lifts to $L$ by D.2.18.
There is an actual cone map
\[
 \mathbb H P^2=S^4\cup_{\nu_4}e^8
              \longrightarrow \cofib(\alpha)
 \tag{D.2.47f}
\]
with top-cell degree one. The pulled-back bottom mod-two class
has an integral lift on $\mathbb H P^2$, so its Pontryagin
square is the mod-four reduction of the integral square and
evaluates to one on the top cell. The defining relative
pairing for $a_Z$ therefore evaluates oddly on $\alpha$.
Since $L$ is six-connected and $H_7(L;\ZZ)=\ZZ/4\{\alpha\}$,
this is the required fundamental class, after fixing its sign.
Changing the chosen nullhomotopy adds the pullback of a class
in $H^7(Y;\ZZ/4)$. Every such class evaluates to zero on
$\alpha$, which factors through $S^4$ and has zero ordinary
Hurewicz image in $Y$. The coefficient universal-coefficient
theorem, using $H_6L=0$, shows that its entire restriction
to $L$ is zero. This proves the stated choice-independence.

On the $K(\FF_2,3)$ fiber of $Z\to Y$, the class $a_Z$
restricts to $p(v_3)=2v_3V_4$ with coefficients $\ZZ/4$.
The coefficient is exactly that of Thomas's formula (2.1),
not a replacement coefficient map. If $v$ is an integral
cochain lifting $v_3$, with $dv=2V$, then
$d(2vV)=4V^2$. Consequently
\[
                  (\beta_4a_Z)|_{K(\FF_2,3)}=V_4^2.
                  \tag{D.2.47g}
\]
Restricting further to $H$ identifies $b_8|_H$.

\medskip\noindent
\textbf{The integral degree-eight extension on $H$.}
We record the required low Eilenberg--MacLane groups, including
their coefficient normalizations:
\[
\begin{array}{c|c|c}
\text{space}&\text{degree}&\text{integral cohomology}\\ \hline
K(\FF_2,3)&8&\ZZ/2\{V_4^2\}\\
K(\FF_2,5)&6,7,8&
 \ZZ/2\{\beta_{\ZZ}s_5\},\ 0,
 \ZZ/2\{\beta_{\ZZ}\operatorname{Sq}^2s_5\}\\
K(\FF_2,4)&9&\ZZ/4\{\beta_4\mathfrak P_2u_4\}\\
K(\FF_2,6)&7,8,9&
 \ZZ/2\{\beta_{\ZZ}w_6\},\ 0,
 \ZZ/2\{\beta_{\ZZ}\operatorname{Sq}^2w_6\}.
\end{array}
\tag{D.2.47h}
\]
These follow directly from Serre's polynomial generators and
the coefficient Bockstein calculation
\cite[\S2, nos.~7--11, pp.~201--205]{I_Serre1953}.
Here are enough details to fix the higher-order entry.
For $K(\FF_2,3)$, the relevant first Bockstein sends
$v_3\operatorname{Sq}^1v_3$ to
$(\operatorname{Sq}^1v_3)^2$; all other degree-eight
classes are either noncycles or first-Bockstein boundaries.
This leaves precisely the displayed order-two integral class.
The next Bockstein page is zero in degrees seven and eight,
so there is no higher-order summand in that degree.
The degree-seven vanishing and the degree-eight generator for
$K(\FF_2,5)$ are in the stable range: the only appropriate
nonzero first Bockstein is
$\operatorname{Sq}^1\operatorname{Sq}^2s_5
=\operatorname{Sq}^3s_5$. The same argument gives the last row.
For $K(\FF_2,4)$ the two surviving first-Bockstein-page classes
in degrees eight and nine are
\[
 u_4^2,\qquad
 u_4\operatorname{Sq}^1u_4+
       \operatorname{Sq}^4\operatorname{Sq}^1u_4.
\]
The Pontryagin-square cochain formula gives the second
Bockstein from the first to the second; equivalently its
integral connecting class is
$\beta_4\mathfrak P_2u_4$, of exact order four. Thus no
additional integral degree-nine summand is omitted.
Explicitly, for an integral lift $U$ with $dU=2A$, use
$U\smile U+U\smile_1dU$ modulo four. Expanding its
coboundary, dividing by four, and reducing modulo two gives
$UA+A\smile_1A$. Its class is
$u_4\operatorname{Sq}^1u_4+
\operatorname{Sq}^4\operatorname{Sq}^1u_4$, the stated
nonzero second-Bockstein target.

Apply the integral Serre sequence to
$K(\FF_2,5)\to H\to K(\FF_2,3)$.
The only pieces in total degree eight are the first and the
degree-eight fiber entries of \textup{(D.2.47h)}.
The base class cannot be hit: its only possible incoming
sources are $(1,6)$ and $(0,7)$, both zero.
The fiber class is universally transgressive, being the
cohomology suspension of
$\beta_{\ZZ}\operatorname{Sq}^2w_6$ in the path fibration
over $K(\FF_2,6)$. Naturality therefore excludes earlier
mixed differentials, including the dimensionally possible
$d_3$ into $(3,6)$. Its transgression is
\[
 \beta_{\ZZ}\operatorname{Sq}^2(v_3^2)
       =\beta_{\ZZ}\bigl((\operatorname{Sq}^1v_3)^2\bigr)=0,
 \tag{D.2.47i}
\]
since the squared class has the integral lift $V_4^2$.
Thus the integral degree-eight group is an extension of
two copies of $\ZZ/2$.

On $H$, Cartan and the nullity of $v_3^2$ give
$\rho_2(V_4^2)=\operatorname{Sq}^2(v_3^2)=0$.
The integral coefficient sequence makes $V_4^2$ divisible
by two. It remains nonzero by the base-class argument above,
so the extension is necessarily $\ZZ/4$ with $V_4^2=2z_8$.
In particular $4z_8=0$, and
$(V_4^2)^2=4z_8^2=0$. This proves \textup{(D.2.47b)}.
One may choose $z_8=\beta_4c_7$, where $c_7$ is a mod-four
lift of $v_3\operatorname{Sq}^1v_3$ on $H$; its first
Bockstein is zero there. This choice is not asserted to
extend to $L$.

\medskip\noindent
\textbf{The mixed differential.}
First use the universal path fibration over $P_6Y$.
The integral sequence for
$K(\FF_2,6)\to P_6Y\to K(\FF_2,4)$ gives
\[
 H^9(P_6Y;\ZZ)=\ZZ/4\{B_9\},
 \qquad B_9=\beta_4\mathfrak P_2u_4.
 \tag{D.2.47j}
\]
Indeed the base degree-nine group in \textup{(D.2.47h)}
has no possible incoming differential. The only other
degree-nine candidate is the universally transgressive fiber
class $\beta_{\ZZ}\operatorname{Sq}^2w_6$. It dies, since
its transgression is
\[
 \beta_{\ZZ}\operatorname{Sq}^2\operatorname{Sq}^3u_4
   =\beta_{\ZZ}\operatorname{Sq}^4\operatorname{Sq}^1u_4\ne0.
 \tag{D.2.47k}
\]
Its mod-two reduction is $(\operatorname{Sq}^1u_4)^2$,
nonzero on the base and unaffected by any earlier incoming
differential. There are no mixed total-degree-nine terms.

Cohomology suspension, the coefficient connecting map, and
Thomas's formula commute. Hence in the universal path fibration
$H\to P(P_6Y)\to P_6Y$ the suspension image is
\[
 \sigma^*B_9=\beta_4p(v_3)=V_4^2=2z_8.
 \tag{D.2.47l}
\]
For a path fibration the transgressive part of fiber cohomology
is exactly this suspension image. Thus $z_8$ must support an
earlier differential. The base is three-connected, and
$H^3(H;\ZZ)=H^5(H;\ZZ)=0$, whereas
$H^4(H;\ZZ)=\ZZ/2\{V_4\}$. The only possible earlier
target is consequently
\[
 E_5^{5,4}=H^5(P_6Y;\ZZ/2)\otimes\{V_4\}
          =\FF_2\{\operatorname{Sq}^1u_4\otimes V_4\}.
\]
No earlier differential hits this group. Necessarily $d_5z_8$
is its nonzero generator.

Finally $H\to L\to Y$ is the pullback of that path fibration
along the actual Postnikov projection. The map carries
$\operatorname{Sq}^1u_4$ to
$\operatorname{Sq}^1U_4=\Sigma^2\operatorname{Sq}^1u_2\ne0$.
Again no earlier differential can hit bidegree $(5,4)$.
Naturality proves \textup{(D.2.47c)}. The square $b_8^2$
has zero restriction to the fiber, but can still be a
nonzero class of positive base filtration; no multiplicative
extension is suppressed in this conclusion.
\end{proof}
